\chapter{การสร้างโมเดิร์นไปป์ไลน์ ELT ด้วย Dataform และ dbt}

การเปลี่ยนกระบวนทัศน์จาก ETL แบบดั้งเดิมสู่ ELT (Extract, Load, Transform) ทำให้การแปลงรูปข้อมูลเกิดขึ้นภายในคลังข้อมูลโดยตรง ซึ่งช่วยเพิ่มความเร็วและลดความซับซ้อนของโครงสร้างพื้นฐาน

\section{สถาปัตยกรรมเหรียญทองสำหรับคลังข้อมูลวิจัย}

ข้อมูลที่ไหลเข้าสู่ระบบถูกจัดระดับความบริสุทธิ์ออกเป็น 3 ชั้น (Medallion Layers):
\begin{enumerate}
    \item \textbf{Bronze Layer (Raw Data):} ข้อมูลดิบจากเซนเซอร์ภาคสนามที่นำเข้าโดยตรง ไม่มีการแปลงสภาพ
    \item \textbf{Silver Layer (Cleaned \& Conformed):} ข้อมูลที่ผ่านการขจัดค่าสูญหาย การปรับสเกล และการตรวจสอบชนิดข้อมูล
    \item \textbf{Gold Layer (Aggregated Business Level):} ตารางสรุปค่าเฉลี่ยสถิติ ดัชนีสุขภาพดิน หรือค่าชี้วัดที่พร้อมสำหรับนำไปวิเคราะห์และสร้างแดชบอร์ด
\end{enumerate}

\section{การสร้าง Dataform Actions ด้วย SQLX}

Dataform ใช้ไฟล์นามสกุล \texttt{.sqlx} ซึ่งผสมผสานระหว่างคำสั่ง SQL และ JavaScript ทำให้สามารถกำหนดตรรกะแบบไดนามิกและการประกาศความสัมพันธ์แบบ Directed Acyclic Graph (DAG) ได้อย่างเป็นระบบ

\begin{activitybox}{การเขียน Assertion เพื่อตรวจสอบคุณภาพข้อมูล}
การกำหนดเงื่อนไข \texttt{assertions} ในไฟล์ SQLX ช่วยให้ระบบหยุดการประมวลผลทันทีหากพบว่าค่าความเป็นกรด-ด่างของดินอยู่นอกช่วงสมเหตุสมผลทางฟิสิกส์ ($0 \le \text{pH} \le 14$)
\end{activitybox}
